\section{Repository}

Questa sezione descrive la repository del progetto QuantSys Institutional, la sua struttura e le modalità di organizzazione del codice e della documentazione.

\subsection{Repository GitHub}

Il progetto QuantSys Institutional è ospitato su una singola repository GitHub:

\begin{itemize}
  \item \textbf{Repository}: \url{https://github.com/Idrorepellente/quantsys_institutional}
\end{itemize}

A differenza di progetti che separano codice e documentazione in repository distinte, QuantSys Institutional adotta una repository monolitica in cui codice, documentazione e file di configurazione coesistono. Questa scelta è motivata dal fatto che molti documenti (es. \texttt{GUIDE\_STRATEGIE\_CUSTOM.md}, \texttt{STRATEGIE\_FX\_CANDIDATE.md}) sono strettamente accoppiati al codice e devono evolvere assieme ad esso.

\subsection{Struttura della repository}

La repository è organizzata in cartelle tematiche, ognuna con una responsabilità ben definita:

\begin{itemize}
  \item \texttt{core/}: tipi base, portfolio, allocator, risk manager, exit policy, indicator pool;
  \item \texttt{backtest/}: engine di backtest event-driven con check intra-bar;
  \item \texttt{data/}: universe, loader yfinance, loader Polygon, loader Dukascopy FX;
  \item \texttt{strategies/}: plugin di strategie atomiche (\texttt{momentum}, \texttt{quality\_lv}, \texttt{value}, \texttt{custom\_example});
  \item \texttt{tests/}: script di walk-forward per ogni strategia, test unitari;
  \item \texttt{scripts/}: script di pipeline (\texttt{build\_feature\_store}, \texttt{run\_ml\_pipeline}, \texttt{bear\_overlay\_backtest}, \texttt{download\_polygon\_m5}, \texttt{download\_dukascopy\_fx});
  \item \texttt{ml/}: pipeline ML con feature engineering e walk-forward training;
  \item \texttt{dashboard/}: Flask app con UI Plotly, API JSON, runner backtest;
  \item \texttt{reporting/}: metriche, factor attribution, ensemble engine;
  \item \texttt{validation/}: walk-forward harness;
  \item \texttt{config/}: file YAML di configurazione (\texttt{quantsys.yaml});
  \item \texttt{output/}: output centralizzati prodotti dai backtest (\texttt{strategies/}, \texttt{ml/});
  \item \texttt{data/cache\_d1/}, \texttt{cache\_polygon\_m5/}, \texttt{cache\_fx\_d1/}, \texttt{cache\_fx\_m5/}: cache di dati di mercato.
\end{itemize}

La radice della repository contiene inoltre i documenti Markdown principali:

\begin{itemize}
  \item \texttt{README.md}: descrizione del progetto;
  \item \texttt{OBIETTIVI.md}: obiettivi misurabili del progetto;
  \item \texttt{GUIDE\_STRATEGIE\_CUSTOM.md}: guida per implementare nuove strategie nel framework plugin;
  \item \texttt{PITCH\_FOF.md}: pitch deck verso Fund of Funds;
  \item \texttt{ROADMAP\_FASE2\_v2.md}: piano per Fase 2 (ampliamento ensemble);
  \item \texttt{STRATEGIE\_FX\_CANDIDATE.md}: proposta di strategie FX da implementare.
\end{itemize}

\subsection{Nomenclatura dei branch}

Il sistema di versionamento prevede la seguente convenzione di naming per i branch:

\begin{itemize}
  \item \texttt{main}: branch principale, sempre in stato deployabile;
  \item \texttt{feature/<nome-strategia>}: per lo sviluppo di una nuova strategia plugin (es. \texttt{feature/stat-arb-coint});
  \item \texttt{feature/<numero-issue>-<descrizione>}: per feature non legate a strategie (es. \texttt{feature/12-exit-policy-engine});
  \item \texttt{fix/<numero-issue>-<descrizione>}: per bug fix (es. \texttt{fix/27-stale-feature-store});
  \item \texttt{docs/<argomento>}: per modifiche esclusivamente documentali (es. \texttt{docs/exit-policy-guide}).
\end{itemize}

Tutti i branch devono essere collegati ad almeno una issue GitHub, che ne motiva l'esistenza e ne traccia lo scopo.

\subsection{Cache dei dati}

Le cartelle \texttt{data/cache\_*} contengono dati di mercato scaricati da provider esterni (Polygon, yfinance, Dukascopy) e \textbf{non vengono versionate su GitHub} (sono escluse via \texttt{.gitignore}). Ogni sviluppatore che clona la repository deve eseguire gli script di download per popolare le cache.

Analogamente, la cartella \texttt{output/} contiene risultati di backtest specifici dell'ambiente di esecuzione e non viene versionata.
